% LaTeX Strategic Analysis for AMC 12B 2017 Problem 16
% Using AMC12/COMC Problem-Solving Framework

\documentclass[]{article}
\usepackage[utf8]{inputenc}
\usepackage{amsmath, amssymb, amsthm}
\usepackage{geometry}
\usepackage{xcolor}
\usepackage[most]{tcolorbox}
\usepackage{enumitem}
\usepackage{fancyhdr}

% Page setup
\geometry{margin=1in}
\pagestyle{fancy}
\fancyhf{}
\rhead{AMC12 Problem Analysis}
\lhead{2017 AMC 12B Problem 16}
\cfoot{\thepage}

% Color definitions
\definecolor{problemconfig}{RGB}{230, 242, 255}
\definecolor{strategic}{RGB}{255, 230, 242}
\definecolor{tactical}{RGB}{242, 255, 230}
\definecolor{techniques}{RGB}{255, 242, 230}
\definecolor{verification}{RGB}{255, 230, 230}
\definecolor{solution}{RGB}{230, 255, 255}
\definecolor{competition}{RGB}{242, 242, 255}
\definecolor{gold}{RGB}{218, 165, 32}

% Box styles
\newtcolorbox{problemconfigbox}{
    colback=problemconfig,
    colframe=blue,
    title=Problem Configuration Analysis,
    fonttitle=\bfseries,
    arc=3pt,
    boxrule=1pt,
    breakable
}

\newtcolorbox{strategicbox}{
    colback=strategic,
    colframe=purple,
    title=Strategic Framework Application,
    fonttitle=\bfseries,
    arc=3pt,
    boxrule=1pt,
    breakable
}

\newtcolorbox{tacticalbox}{
    colback=tactical,
    colframe=green,
    title=Tactical Methodology Selection,
    fonttitle=\bfseries,
    arc=3pt,
    boxrule=1pt,
    breakable
}

\newtcolorbox{techniquesbox}{
    colback=techniques,
    colframe=orange,
    title=Conversion Techniques Application,
    fonttitle=\bfseries,
    arc=3pt,
    boxrule=1pt,
    breakable
}

\newtcolorbox{verificationbox}{
    colback=verification,
    colframe=red,
    title=Verification Strategy Design,
    fonttitle=\bfseries,
    arc=3pt,
    boxrule=1pt,
    breakable
}

\newtcolorbox{solutionbox}{
    colback=solution,
    colframe=cyan,
    title=Solution Roadmap,
    fonttitle=\bfseries,
    arc=3pt,
    boxrule=1pt,
    breakable
}

\newtcolorbox{competitionbox}{
    colback=competition,
    colframe=purple,
    title=Competition Strategy Integration,
    fonttitle=\bfseries,
    arc=3pt,
    boxrule=1pt,
    breakable
}

\newtcolorbox{keyinsightbox}{
    colback=yellow!20,
    colframe=gold,
    title=Key Insight,
    fonttitle=\bfseries,
    arc=3pt,
    boxrule=2pt,
    leftrule=5pt,
    breakable
}

\newtcolorbox{warningbox}{
    colback=red!10,
    colframe=red!50,
    title=Warning: Common Pitfall,
    fonttitle=\bfseries,
    arc=3pt,
    boxrule=2pt,
    leftrule=5pt,
    breakable
}

\newtcolorbox{tipbox}{
    colback=green!10,
    colframe=green!50,
    title=Pro Tip,
    fonttitle=\bfseries,
    arc=3pt,
    boxrule=2pt,
    leftrule=5pt,
    breakable
}

\begin{document}

\title{\textbf{AMC12/COMC Strategic Problem Analysis\\2017 AMC 12B Problem 16}}
\author{Competition Math Framework}
\date{\today}
\maketitle

\section*{Problem Statement}
The number $21! = 51,090,942,171,709,440,000$ has over $60,000$ positive integer divisors. One of them is chosen at random. What is the probability that it is odd?

\textbf{Answer Choices:}
\[\textbf{(A)}\ \frac{1}{21} \qquad \textbf{(B)}\ \frac{1}{19} \qquad \textbf{(C)}\ \frac{1}{18} \qquad \textbf{(D)}\ \frac{1}{2} \qquad \textbf{(E)}\ \frac{11}{21}\]

\begin{problemconfigbox}
\subsection*{A. Problem Classification}
\begin{itemize}[leftmargin=*]
    \item \textbf{Problem Type}: To Find (compute a probability)
    \item \textbf{Difficulty Band}: Mid-band (Problem 16 of 25) - requires understanding of prime factorization and divisor counting
    \item \textbf{Competition Context}: AMC12 (75 minutes, 25 problems) - approximately 3 minutes allocated
    \item \textbf{Mathematical Domain}: Number Theory with Probability elements (divisor properties, prime factorization, counting)
\end{itemize}

\subsection*{B. Problem Structure Identification}
\begin{itemize}[leftmargin=*]
    \item \textbf{Given Information}: 
    \begin{itemize}
        \item $21! = 51,090,942,171,709,440,000$
        \item $21!$ has over 60,000 positive integer divisors
        \item One divisor is chosen at random
    \end{itemize}
    \item \textbf{Constraints}: 
    \begin{itemize}
        \item Must work with factorial $21!$
        \item Random selection means uniform probability distribution
        \item Divisor must be positive integer
    \end{itemize}
    \item \textbf{Objective}: Find the probability that a randomly chosen divisor is odd
    \item \textbf{Hidden Assumptions}: 
    \begin{itemize}
        \item A number is odd if and only if it contains no factors of 2
        \item Prime factorization uniquely determines divisors
        \item All divisors are equally likely to be chosen
    \end{itemize}
    \item \textbf{Answer Choice Leverage}: 
    \begin{itemize}
        \item Answer (D) $\frac{1}{2}$ would be the naive guess (half odd, half even)
        \item Small denominators suggest we need to count something specific
        \item Denominator 19 in choice (B) hints at $2^{18}$ (since $18 + 1 = 19$)
    \end{itemize}
\end{itemize}

\subsection*{C. Strategic Orientation}
\begin{itemize}[leftmargin=*]
    \item \textbf{Hypothesis}: A divisor of $21!$ is determined by choosing exponents in the prime factorization
    \item \textbf{Conclusion}: Need to find the ratio of odd divisors to total divisors
    \item \textbf{Notation}: 
    \begin{itemize}
        \item $21! = 2^{e_2} \cdot 3^{e_3} \cdot 5^{e_5} \cdot 7^{e_7} \cdot 11^{e_{11}} \cdot 13^{e_{13}} \cdot 17^{e_{17}} \cdot 19^{e_{19}}$
        \item $e_2$ = exponent of 2 in prime factorization
    \end{itemize}
    \item \textbf{Intuitive Approach}: 
    \begin{itemize}
        \item An odd number has no factors of 2
        \item Count divisors with zero factors of 2 versus all divisors
        \item Use Legendre's formula to find $e_2$
    \end{itemize}
    \item \textbf{Cross-Domain Potential}: Could use combinatorics (choosing exponents) or direct probability computation
\end{itemize}
\end{problemconfigbox}

\begin{strategicbox}
\subsection*{A. Psychological Strategies}
\begin{itemize}[leftmargin=*]
    \item \textbf{Mental Toughness}: If Legendre's formula is unfamiliar, use systematic counting: $\lfloor 21/2 \rfloor + \lfloor 21/4 \rfloor + \lfloor 21/8 \rfloor + \lfloor 21/16 \rfloor$. This is a mechanical process that builds confidence.
    \item \textbf{Creativity Cultivation}: Consider the problem from multiple angles - divisor counting, probability theory, or combinatorial selection. Each perspective reinforces the others.
    \item \textbf{Iterative Refinement}: If the first attempt yields an answer not in the choices, double-check the exponent calculation. This is a common error point.
\end{itemize}

\subsection*{B. Investigation Strategies}
\begin{itemize}[leftmargin=*]
    \item \textbf{Startup Orientation}: Recognize this as a number theory divisor problem. Prime factorization is key.
    \item \textbf{Get Your Hands Dirty}: Test with smaller factorial like $5! = 120 = 2^3 \cdot 3 \cdot 5$. Divisors: 1, 2, 3, 4, 5, 6, 8, 10, 12, 15, 20, 24, 30, 40, 60, 120 (16 total). Odd ones: 1, 3, 5, 15 (4 total). Probability = $\frac{4}{16} = \frac{1}{4} = \frac{1}{3+1}$. Pattern emerges!
    \item \textbf{Penultimate Step}: Once we know $e_2 = 18$, the final step is recognizing that probability = $\frac{1}{e_2 + 1} = \frac{1}{19}$
    \item \textbf{Wishful Thinking}: Wish we could separate the ``even'' part from the ``odd'' part of $21!$. This is exactly what prime factorization accomplishes!
    \item \textbf{Make It Easier}: Try $4! = 24 = 2^3 \cdot 3$. Total divisors: $(3+1)(1+1) = 8$. Odd divisors (no factor of 2): $(0+1)(1+1) = 2$ (namely 1 and 3). Probability = $\frac{2}{8} = \frac{1}{4}$, confirming pattern.
    \item \textbf{Conjecture via Small Cases}: Pattern from $4!$ and $5!$: probability = $\frac{1}{\text{exponent of 2} + 1}$
    \item \textbf{Visualization}: Not directly visual, but can think of a tree structure where each branch represents choosing an exponent for a prime
\end{itemize}

\subsection*{C. Late-Band Strategic Playbook}
\begin{itemize}[leftmargin=*]
    \item \textbf{Answer-Choice Leverage}: Eliminate (D) immediately - factorials have far more even divisors than odd ones. Denominator 19 in (B) strongly suggests $2^{18}$.
    \item \textbf{Make It Easier}: Work with $5!$ or $10!$ to understand the pattern before tackling $21!$
    \item \textbf{Penultimate Step}: The final calculation after finding $e_2 = 18$ is trivial: $\frac{1}{19}$
    \item \textbf{Wishful Thinking}: If only we could count odd divisors separately from even ones... which is exactly what happens!
    \item \textbf{Conjecture via Small Cases}: Test formula with $3! = 6 = 2^1 \cdot 3$. Predicted: $\frac{1}{2}$. Actual: divisors are 1,2,3,6; odd are 1,3; probability = $\frac{2}{4} = \frac{1}{2}$. $\checkmark$
\end{itemize}

\subsection*{D. Argument Construction Strategy}
\begin{itemize}[leftmargin=*]
    \item \textbf{Direct Proof}: Use divisor counting formula and Legendre's formula directly
    \item \textbf{Proof by Contradiction}: Not applicable here
    \item \textbf{Mathematical Induction}: Could prove Legendre's formula inductively, but not needed for this problem
    \item \textbf{Stylistic Conventions}: Clear statement that odd divisors correspond to choosing exponent 0 for prime 2
\end{itemize}

\subsection*{E. Cross-Domain Translation}
\begin{itemize}[leftmargin=*]
    \item \textbf{Draw a Picture}: Can visualize as a lattice of exponent choices
    \item \textbf{Recast the Problem}: From probability to pure counting: ``How many divisors of $21!$ are odd?''
    \item \textbf{Change Your Point of View}: Think of divisors as products of prime powers. Odd divisors exclude the prime 2.
\end{itemize}
\end{strategicbox}

\begin{tacticalbox}
\subsection*{A. Core Mathematical Tactics}
\begin{itemize}[leftmargin=*]
    \item \textbf{Symmetry}: Not directly applicable
    \item \textbf{The Extreme Principle}: Not directly applicable
    \item \textbf{The Pigeonhole Principle}: Not directly applicable
    \item \textbf{Invariants}: The prime factorization structure is invariant - key to the solution
\end{itemize}

\subsection*{B. Crossover Tactics}
\begin{itemize}[leftmargin=*]
    \item \textbf{Graph Theory}: Not applicable
    \item \textbf{Complex Numbers}: Not applicable
    \item \textbf{Generating Functions}: Could use generating functions for divisor counting, but overkill here
\end{itemize}
\end{tacticalbox}

\begin{techniquesbox}
\subsection*{A. Recognition Index - Instant Identification}
\begin{itemize}[leftmargin=*]
    \item \textbf{Trigger Pattern Detected}: ``Divisors of $n!$'' + ``odd divisor'' → Prime factorization + Legendre's formula
    \item \textbf{Domain Signal}: Number Theory - divisor counting
    \item \textbf{Structural Signature}: Factorial with probability question about divisor property
\end{itemize}

\subsection*{B. High-Probability Techniques (Selected)}
\begin{itemize}[leftmargin=*]
    \item \textbf{Prime Factorization Analysis}: $\checkmark$ CRITICAL - Must find $21! = 2^{18} \cdot 3^9 \cdot 5^4 \cdot 7^3 \cdot 11 \cdot 13 \cdot 17 \cdot 19$
    \item \textbf{Legendre's Formula}: $\checkmark$ ESSENTIAL - $e_p(n!) = \sum_{i=1}^{\infty} \lfloor n/p^i \rfloor$
    \item \textbf{Divisor Counting Formula}: $\checkmark$ REQUIRED - Total divisors = $(e_2+1)(e_3+1) \cdots (e_{19}+1)$
    \item \textbf{Fundamental Counting Principle}: $\checkmark$ NEEDED - Each prime's exponent chosen independently
\end{itemize}

\subsection*{C. Technique Selection Rationale}
\begin{itemize}[leftmargin=*]
    \item \textbf{Primary Technique}: Legendre's Formula + Divisor Counting
    \item \textbf{Backup Techniques}: Manual counting of multiples (slower but reliable)
    \item \textbf{Technique Integration}: Combine Legendre for finding $e_2$, then use divisor formula for counting
    \item \textbf{Conflict Resolution}: No conflicts - the techniques complement each other perfectly
\end{itemize}

\subsection*{D. Detailed Technique Application}

\textbf{Step 1: Apply Legendre's Formula for $e_2$}

Find the exponent of prime $p=2$ in $21!$:
\begin{align*}
e_2(21!) &= \lfloor 21/2 \rfloor + \lfloor 21/4 \rfloor + \lfloor 21/8 \rfloor + \lfloor 21/16 \rfloor + \lfloor 21/32 \rfloor + \cdots\\
&= 10 + 5 + 2 + 1 + 0 + \cdots\\
&= 18
\end{align*}

\textbf{Step 2: Understand Divisor Structure}

Any divisor of $21!$ has the form:
\[d = 2^{a_2} \cdot 3^{a_3} \cdot 5^{a_5} \cdot 7^{a_7} \cdot 11^{a_{11}} \cdot 13^{a_{13}} \cdot 17^{a_{17}} \cdot 19^{a_{19}}\]
where $0 \leq a_2 \leq 18$, $0 \leq a_3 \leq e_3$, etc.

\textbf{Step 3: Count Total Divisors}

By the Fundamental Counting Principle:
\[\text{Total divisors} = (18+1) \cdot (e_3+1) \cdot (e_5+1) \cdot \cdots \cdot (e_{19}+1)\]

\textbf{Step 4: Count Odd Divisors}

A divisor is odd if and only if $a_2 = 0$. Thus:
\[\text{Odd divisors} = (1) \cdot (e_3+1) \cdot (e_5+1) \cdot \cdots \cdot (e_{19}+1)\]

\textbf{Step 5: Calculate Probability}

\begin{align*}
P(\text{odd}) &= \frac{\text{Odd divisors}}{\text{Total divisors}}\\
&= \frac{(1) \cdot (e_3+1) \cdot (e_5+1) \cdots (e_{19}+1)}{(19) \cdot (e_3+1) \cdot (e_5+1) \cdots (e_{19}+1)}\\
&= \frac{1}{19}
\end{align*}

Notice the beautiful cancellation - we don't even need to compute the other exponents!

\subsection*{E. Integration Strategy}
The techniques integrate seamlessly:
\begin{enumerate}
    \item Legendre's formula identifies $e_2 = 18$
    \item Divisor counting formula structures the problem
    \item Fundamental counting principle enables the calculation
    \item Cancellation simplifies to the final answer
\end{enumerate}
\end{techniquesbox}

\begin{verificationbox}
\subsection*{A. Verification Level Selection}
\begin{itemize}[leftmargin=*]
    \item \textbf{Problem Difficulty}: Mid-band (P16)
    \item \textbf{Confidence Level}: High (90\%+) after calculation
    \item \textbf{Available Time}: 30-45 seconds for verification
    \item \textbf{Cost of Errors}: Moderate - but P16 should be solved correctly
    \item \textbf{Selected Level}: Level 4 (Standard Verification)
\end{itemize}

\subsection*{B. Standard Verification Methods Applied}

\textbf{Primary Verification - Recalculate Legendre:}
\begin{itemize}[leftmargin=*]
    \item $\lfloor 21/2 \rfloor = 10$ $\checkmark$
    \item $\lfloor 21/4 \rfloor = 5$ $\checkmark$
    \item $\lfloor 21/8 \rfloor = 2$ $\checkmark$
    \item $\lfloor 21/16 \rfloor = 1$ $\checkmark$
    \item Sum: $10 + 5 + 2 + 1 = 18$ $\checkmark$
\end{itemize}

\textbf{Boundary Check:}
\begin{itemize}[leftmargin=*]
    \item For $n=1$: divisor is 1 (odd), probability = 1. Formula: undefined (no factors of 2). Makes sense!
    \item For $n=2$: $2! = 2$, divisors: 1, 2. Odd: just 1. Probability = $\frac{1}{2} = \frac{1}{1+1}$ $\checkmark$
\end{itemize}

\textbf{Alternative Approach Glimpse - Small Case Verification:}

Test with $5! = 120 = 2^3 \cdot 3 \cdot 5$:
\begin{itemize}[leftmargin=*]
    \item Predicted probability: $\frac{1}{3+1} = \frac{1}{4}$
    \item Total divisors: $(3+1)(1+1)(1+1) = 16$
    \item Odd divisors: $(1)(1+1)(1+1) = 4$ (namely: 1, 3, 5, 15)
    \item Actual probability: $\frac{4}{16} = \frac{1}{4}$ $\checkmark$
\end{itemize}

\subsection*{C. Domain-Specific Verification}

\textbf{Number Theory Consistency:}
\begin{itemize}[leftmargin=*]
    \item An odd number has $2^0$ in its factorization: $\checkmark$
    \item The exponent of 2 in $21!$ should be less than 21: $18 < 21$ $\checkmark$
    \item Answer $\frac{1}{19}$ is much less than $\frac{1}{2}$, reflecting that most divisors of factorials are even: $\checkmark$
\end{itemize}

\textbf{Answer Choice Consistency:}
\begin{itemize}[leftmargin=*]
    \item Our answer $\frac{1}{19}$ matches choice (B)
    \item Denominator 19 = $18 + 1$ makes sense given $e_2 = 18$
    \item Value is between 0 and 1 as required for probability
\end{itemize}

\subsection*{D. Confidence Calibration}
\begin{itemize}[leftmargin=*]
    \item \textbf{Initial Confidence}: 85\% (clear method but calculation-dependent)
    \item \textbf{Post-Verification}: 95\% (multiple checks confirm answer)
    \item \textbf{Remaining Uncertainty}: 5\% (possible arithmetic error in Legendre's formula)
    \item \textbf{Decision}: Mark answer (B), move on with high confidence
\end{itemize}
\end{verificationbox}

\begin{solutionbox}
\subsection*{A. Complete Solution Path}

\textbf{Step 1: Understand the Problem Structure (30 seconds)}
\begin{itemize}[leftmargin=*]
    \item Need probability that random divisor of $21!$ is odd
    \item Odd means no factors of 2
    \item Will use prime factorization and divisor counting
\end{itemize}

\textbf{Step 2: Find Exponent of 2 in $21!$ (60 seconds)}

Using Legendre's Formula:
\[e_2(21!) = \lfloor 21/2 \rfloor + \lfloor 21/4 \rfloor + \lfloor 21/8 \rfloor + \lfloor 21/16 \rfloor = 10 + 5 + 2 + 1 = 18\]

Therefore $21! = 2^{18} \cdot (\text{odd number})$

\textbf{Step 3: Set Up Divisor Counting (30 seconds)}

Any divisor has form $d = 2^{a_2} \cdot (\text{odd divisor})$ where $0 \leq a_2 \leq 18$

Total divisors: Each choice of $a_2$ paired with each choice of odd divisor

Total ways to choose $a_2$: $19$ choices (from 0 to 18)

\textbf{Step 4: Count Odd Divisors (20 seconds)}

Odd divisors correspond to $a_2 = 0$ only

Number of odd divisors: 1 choice for $a_2$ times (number of odd divisors of odd part)

\textbf{Step 5: Calculate Probability (20 seconds)}

\[P(\text{odd}) = \frac{\text{\# odd divisors}}{\text{\# total divisors}} = \frac{1 \cdot N}{19 \cdot N} = \frac{1}{19}\]

where $N$ is the number of divisors of the odd part (cancels out!)

\textbf{Step 6: Verify and Select Answer (20 seconds)}

Answer: $\boxed{\textbf{(B)}\ \frac{1}{19}}$

\subsection*{B. Alternative Solution Approaches}

\textbf{Alternative 1 - Complete Prime Factorization:}

Find all exponents, compute total and odd divisors separately:
\begin{itemize}[leftmargin=*]
    \item $21! = 2^{18} \cdot 3^9 \cdot 5^4 \cdot 7^3 \cdot 11 \cdot 13 \cdot 17 \cdot 19$
    \item Total divisors: $19 \cdot 10 \cdot 5 \cdot 4 \cdot 2 \cdot 2 \cdot 2 \cdot 2 = 60,800$
    \item Odd divisors: $1 \cdot 10 \cdot 5 \cdot 4 \cdot 2 \cdot 2 \cdot 2 \cdot 2 = 3,200$
    \item Probability: $\frac{3200}{60800} = \frac{1}{19}$ $\checkmark$
\end{itemize}

(Much more work but gives same answer!)

\textbf{Alternative 2 - Pattern Recognition from Small Cases:}

Compute for small factorials, notice pattern, apply to $21!$

\subsection*{C. Pitfall Avoidance Strategy}

\textbf{Common Error 1: Miscounting powers of 2}
\begin{itemize}[leftmargin=*]
    \item Prevention: Write out all terms in Legendre's formula explicitly
    \item Check: $\lfloor 21/32 \rfloor = 0$, so stop at $2^{16}$
\end{itemize}

\textbf{Common Error 2: Confusing number of choices with exponent}
\begin{itemize}[leftmargin=*]
    \item Prevention: Remember $2^{18}$ gives 19 choices (0 through 18)
    \item Mnemonic: ``exponent + 1 = choices''
\end{itemize}

\textbf{Common Error 3: Thinking probability should be close to $\frac{1}{2}$}
\begin{itemize}[leftmargin=*]
    \item Prevention: Factorials have many more even divisors than odd ones
    \item Intuition: Every even number 2, 4, 6, ..., 20 contributes factors of 2
\end{itemize}

\textbf{Common Error 4: Not recognizing cancellation}
\begin{itemize}[leftmargin=*]
    \item Prevention: Set up ratio carefully, notice odd part cancels
    \item Benefit: Saves time computing other exponents
\end{itemize}
\end{solutionbox}

\begin{competitionbox}
\subsection*{A. Time Management}
\begin{itemize}[leftmargin=*]
    \item \textbf{Time Budget}: 3 minutes for P16 (mid-band problem)
    \item \textbf{Actual Expected}: 2.5-3 minutes with verification
    \item \textbf{Decision Point}: If not solved in 4 minutes, flag and move on
    \item \textbf{Time Distribution}:
    \begin{itemize}
        \item Problem reading and setup: 30 seconds
        \item Legendre's formula calculation: 60 seconds
        \item Divisor analysis: 50 seconds
        \item Verification: 30 seconds
    \end{itemize}
\end{itemize}

\subsection*{B. Problem Prioritization Strategy}
\begin{itemize}[leftmargin=*]
    \item \textbf{Accessibility}: HIGH - standard technique with clear path
    \item \textbf{Point Value}: Standard (6 points for AMC12)
    \item \textbf{Strategic Importance}: ``Must solve'' for scores above 120
    \item \textbf{Domain Rotation}: Good balance - number theory after geometry/algebra problems
\end{itemize}

\subsection*{C. Risk Management}

\textbf{Confidence Thresholds:}
\begin{itemize}[leftmargin=*]
    \item \textbf{Target Confidence}: 90\%+ (achievable with verification)
    \item \textbf{Error Probability}: Low if Legendre's formula is applied correctly
    \item \textbf{Review Priority}: Low - move forward confidently if verified
\end{itemize}

\textbf{Answer Strategy:}
\begin{itemize}[leftmargin=*]
    \item Eliminate (D) immediately (naive $\frac{1}{2}$ is wrong)
    \item If uncertain after calculation, choice (B) has denominator matching pattern
    \item Blank vs. Guess: If completely stuck, educated guess favors (B) based on $e_2 + 1$ pattern
\end{itemize}

\subsection*{D. Problem Abandonment Criteria}

\textbf{Soft Abandon} (after 4 minutes):
\begin{itemize}[leftmargin=*]
    \item If Legendre's formula is unfamiliar and manual counting is taking too long
    \item Mark answer (B) based on pattern recognition
    \item Flag for review if time permits
\end{itemize}

\textbf{Hard Abandon} (immediate):
\begin{itemize}[leftmargin=*]
    \item \textit{Not applicable} - this problem has clear technique pathway
    \item If completely unfamiliar with prime factorization, skip to P17
\end{itemize}

\textbf{Optimization Strategy - AMC12 Specific:}
\begin{itemize}[leftmargin=*]
    \item Blanks worth 1.5 points, wrong answers worth 0
    \item With solid approach, wrong answer here unlikely
    \item If rushed, use answer choice leverage: (B) has special structure
\end{itemize}
\end{competitionbox}

\begin{keyinsightbox}
\textbf{Key Insight}: The beautiful cancellation is the heart of this problem! 

We don't need to compute exponents for primes 3, 5, 7, etc. because:
\[\frac{\text{Odd divisors}}{\text{Total divisors}} = \frac{1 \cdot (\text{divisors of odd part})}{19 \cdot (\text{divisors of odd part})} = \frac{1}{19}\]

The ``divisors of odd part'' appears in both numerator and denominator, so it cancels completely. This is why Legendre's formula for $e_2$ is the \textit{only} calculation needed.

\textbf{Pattern Recognition}: For any factorial $n!$ with $2^k$ as highest power of 2 dividing it:
\[P(\text{divisor is odd}) = \frac{1}{k+1}\]
\end{keyinsightbox}

\begin{warningbox}
\textbf{Warning}: The most common error is miscounting the exponent of 2!

Students often forget one term in Legendre's formula or make arithmetic mistakes:
\begin{itemize}[leftmargin=*]
    \item Double-check: $10 + 5 + 2 + 1 = 18$ (not 16, not 19)
    \item Remember to stop when $\lfloor 21/2^k \rfloor = 0$
    \item The next term would be $\lfloor 21/32 \rfloor = 0$, contributing nothing
\end{itemize}

Also watch out for the off-by-one error: $2^{18}$ gives 19 possible exponents (0 through 18), not 18!
\end{warningbox}

\begin{tipbox}
\textbf{Pro Tip - Competition Time Saver}: 

If you know Legendre's formula cold, you can write the answer in under 90 seconds:
\begin{enumerate}
    \item Compute $e_2(21!) = 18$ (45 sec)
    \item Recognize cancellation pattern (15 sec)
    \item Write answer $\frac{1}{19}$ (10 sec)
    \item Quick verification (20 sec)
\end{enumerate}

\textbf{Practice Recommendation}: Memorize Legendre's formula and practice computing it for factorials $10!$, $15!$, $20!$, $25!$, $30!$ to build speed.

\textbf{Answer Choice Hack}: In AMC probability problems involving factorials and divisors, if you see denominators like 17, 19, 23 (primes or near-primes), they often equal (exponent + 1) for some prime's exponent in the factorial.
\end{tipbox}

\section*{Final Answer}
\[\boxed{\textbf{(B)}\ \frac{1}{19}}\]

\section*{Confidence Assessment}
\begin{itemize}[leftmargin=*]
    \item \textbf{Confidence Level}: 95\% - standard technique properly applied and verified
    \item \textbf{Verification Applied}: Level 4 Standard Verification (Legendre recalculation, small case check, answer choice consistency)
    \item \textbf{Flag for Review}: No - high confidence, move forward
    \item \textbf{Time Spent}: 2.5 minutes (on target for P16)
    \item \textbf{Key Success Factors}: 
    \begin{itemize}
        \item Immediate recognition of Legendre's formula application
        \item Clean cancellation reducing computation
        \item Multiple verification checks confirming answer
    \end{itemize}
\end{itemize}

\end{document}